\documentclass[12pt,a4paper]{article}

\usepackage[T1]{fontenc}
\usepackage[utf8]{inputenc}
\usepackage[romanian]{babel}
\usepackage{geometry}
\usepackage{amsmath}

\geometry{margin=2.5cm}

\title{Tutoriat 3\\Cautare Adversariala - Introducere in cautare adversariala}
\author{}
\date{}

\begin{document}

\maketitle


\section{Introducere}

Inteligența Artificială (Artificial Intelligence - AI) reprezintă domeniul informaticii care studiază dezvoltarea sistemelor capabile să realizeze sarcini ce necesită în mod obișnuit inteligență umană.

Printre obiectivele principale ale Inteligenței Artificiale se numără:

\begin{itemize}
\item rezolvarea de probleme;
\item luarea deciziilor;
\item planificarea acțiunilor;
\item învățarea din experiență;
\item înțelegerea limbajului natural;
\item percepția mediului înconjurător.
\end{itemize}

Un concept central în Inteligența Artificială este cel de agent inteligent.

Un agent este o entitate care percepe mediul prin intermediul senzorilor și acționează asupra acestuia folosind actuatori.

Exemple de agenți:

\begin{itemize}
\item un robot autonom;
\item un program care joacă șah;
\item un vehicul autonom;
\item un asistent virtual;
\item un sistem de recomandare.
\end{itemize}

Pentru ca un agent să poată lua decizii eficiente, acesta trebuie să analizeze mediul în care operează și să aleagă acțiunile care conduc la obținerea unui rezultat cât mai favorabil.

În cadrul acestui tutoriat vom studia o categorie importantă de probleme din Inteligența Artificială: problemele de căutare adversarială.

Aceste probleme apar în jocuri și scenarii multi-agent în care succesul unui agent depinde de acțiunile altor agenți.

Exemple clasice:

\begin{itemize}
\item șah;
\item dame;
\item X și 0;
\item Go;
\item Poker.
\end{itemize}

Pentru rezolvarea acestor probleme vom introduce:

\begin{itemize}
\item arborii de joc;
\item algoritmul MiniMax;
\item funcțiile de evaluare;
\item tehnica Alpha-Beta Pruning.
\end{itemize}

Aceste metode permit unui agent să selecteze mutarea optimă presupunând că adversarul joacă în mod rațional și încearcă să îi minimizeze șansele de succes.












\section{Agenți inteligenți și medii}

Un concept fundamental în Inteligența Artificială este cel de agent inteligent.

Un agent este o entitate care percepe mediul în care operează și acționează asupra acestuia astfel încât să își atingă obiectivele.

Schema generală este:

\begin{center}
Mediu → Percepție → Agent → Acțiune → Mediu
\end{center}

Agentul primește informații despre mediu, le analizează și selectează o acțiune pe care o consideră optimă.

\subsection{Percepții și acțiuni}

Percepțiile reprezintă informațiile pe care agentul le obține din mediu.

Acțiunile reprezintă operațiile pe care agentul le poate executa pentru a modifica mediul sau pentru a influența evoluția acestuia.

Exemple:

\begin{itemize}
\item un robot percepe distanțe și se poate deplasa;
\item un vehicul autonom percepe obstacole și poate accelera sau frâna;
\item un program de șah percepe configurația tablei și poate efectua mutări.
\end{itemize}

\subsection{Raționalitatea unui agent}

Un agent este considerat rațional dacă selectează acțiunea care maximizează performanța sa pe baza informațiilor disponibile.

Agentul nu trebuie să fie perfect.

El trebuie să ia cea mai bună decizie posibilă folosind cunoștințele și resursele de care dispune la momentul respectiv.

\subsection{Mediul}

Mediul reprezintă tot ceea ce se află în afara agentului și cu care acesta interacționează.

Caracteristicile mediului influențează dificultatea problemei pe care agentul trebuie să o rezolve.

\subsection{Tipuri de medii}

\subsubsection{Complet observabil vs parțial observabil}

Un mediu este complet observabil dacă agentul poate cunoaște în orice moment întreaga stare a sistemului.

Exemple:

\begin{itemize}
\item șah;
\item X și 0.
\end{itemize}

Un mediu este parțial observabil dacă anumite informații lipsesc sau sunt inaccesibile.

Exemple:

\begin{itemize}
\item poker;
\item conducerea unui autovehicul în trafic.
\end{itemize}

\subsubsection{Determinist vs stochastic}

Într-un mediu determinist, aceeași acțiune produce întotdeauna același rezultat.

Exemplu:

\begin{itemize}
\item șah.
\end{itemize}

Într-un mediu stochastic există elemente aleatorii care influențează rezultatul.

Exemple:

\begin{itemize}
\item zaruri;
\item jocuri de noroc;
\item vremea.
\end{itemize}

\subsubsection{Static vs dinamic}

Un mediu static nu se modifică în timp ce agentul analizează situația.

Exemplu:

\begin{itemize}
\item rezolvarea unui puzzle.
\end{itemize}

Un mediu dinamic se poate modifica independent de agent.

Exemple:

\begin{itemize}
\item traficul rutier;
\item jocurile în timp real.
\end{itemize}

\subsubsection{Discret vs continuu}

Un mediu discret are un număr finit de stări și acțiuni.

Exemple:

\begin{itemize}
\item șah;
\item dame;
\item X și 0.
\end{itemize}

Un mediu continuu permite valori reale pentru stări sau acțiuni.

Exemple:

\begin{itemize}
\item controlul unui robot;
\item pilotarea unui avion.
\end{itemize}

\subsection{Agenți în jocuri}

În cadrul jocurilor, fiecare jucător poate fi modelat ca un agent.

Fiecare agent încearcă să își maximizeze propriul succes și să reducă șansele adversarului.

Această situație conduce la probleme de căutare adversarială, unde decizia optimă depinde atât de acțiunile proprii, cât și de reacțiile adversarului.

Studiul acestor probleme va conduce la introducerea arborilor de joc și a algoritmului MiniMax.








\section{Căutare adversarială}

În multe aplicații de Inteligență Artificială, agentul nu operează singur în mediu, ci trebuie să ia decizii în prezența altor agenți care urmăresc obiective diferite sau chiar opuse.

Aceste situații conduc la probleme de căutare adversarială.

Căutarea adversarială studiază metodele prin care un agent poate determina cea mai bună acțiune ținând cont de faptul că adversarul va încerca să îi reducă șansele de succes.

\subsection{Motivație}

În problemele clasice de căutare, mediul este pasiv.

Agentul trebuie doar să găsească o secvență de acțiuni care conduce la atingerea scopului.

În jocuri însă, situația este diferită.

După fiecare decizie a agentului, adversarul poate răspunde cu propria mutare.

Astfel, rezultatul final depinde de deciziile ambilor jucători.

\subsection{Exemple}

Problemele de căutare adversarială apar în:

\begin{itemize}
\item șah;
\item dame;
\item X și 0;
\item Go;
\item Reversi;
\item poker;
\item jocuri strategice.
\end{itemize}

În toate aceste exemple, succesul unui jucător depinde atât de propriile mutări, cât și de mutările adversarului.

\subsection{Modelarea unei probleme de joc}

Pentru a putea analiza un joc, trebuie definite următoarele elemente:

\begin{itemize}
\item starea inițială;
\item mulțimea acțiunilor posibile;
\item funcția de tranziție;
\item stările terminale;
\item funcția de utilitate.
\end{itemize}

Aceste componente permit construirea unui model formal al jocului.

\subsection{Stări}

O stare descrie complet configurația curentă a jocului.

Exemple:

\begin{itemize}
\item poziția pieselor pe tabla de șah;
\item configurația tablei de X și 0;
\item distribuția cărților într-un joc de poker.
\end{itemize}

Din fiecare stare se pot genera stările succesoare prin aplicarea mutărilor legale.

\subsection{Stări terminale}

O stare terminală reprezintă o stare în care jocul s-a încheiat.

Exemple:

\begin{itemize}
\item șah mat;
\item remiză;
\item victoria unuia dintre jucători;
\item ocuparea completă a tablei la X și 0.
\end{itemize}

Într-o stare terminală nu mai există mutări posibile.

\subsection{Funcția de utilitate}

Funcția de utilitate atribuie o valoare numerică fiecărei stări terminale.

Această valoare exprimă cât de favorabil este rezultatul pentru jucătorul analizat.

Un exemplu simplu:

\begin{itemize}
\item +1 pentru victorie;
\item 0 pentru remiză;
\item -1 pentru înfrângere.
\end{itemize}

Scopul agentului este maximizarea utilității.

\subsection{Jucătorul MAX și jucătorul MIN}

În teoria jocurilor, participanții sunt reprezentați de obicei astfel:

\begin{itemize}
\item MAX - jucătorul pentru care calculăm strategia optimă;
\item MIN - adversarul, care încearcă să minimizeze câștigul lui MAX.
\end{itemize}

La fiecare nivel al arborelui de joc, unul dintre cei doi jucători alege mutarea următoare.

\subsection{Obiectivul căutării adversariale}

Problema principală este:

\begin{center}
"Ce mutare trebuie să aleagă MAX dacă presupunem că MIN joacă perfect?"
\end{center}

Pentru a răspunde la această întrebare este necesară generarea și analiza arborelui de joc.

Acest proces conduce la algoritmul MiniMax, una dintre cele mai importante metode de căutare adversarială din Inteligența Artificială.








\section{Jocuri și tipuri de jocuri}

Jocurile reprezintă una dintre cele mai importante aplicații ale căutării adversariale în Inteligența Artificială.

Într-un joc există mai mulți agenți care iau decizii succesive pentru atingerea propriilor obiective.

Succesul unui jucător depinde atât de propriile mutări, cât și de deciziile adversarilor.

\subsection{Componentele unui joc}

Un joc poate fi descris prin următoarele elemente:

\begin{itemize}
\item mulțimea jucătorilor;
\item starea inițială;
\item acțiunile disponibile;
\item regulile de tranziție între stări;
\item stările terminale;
\item funcția de utilitate.
\end{itemize}

Aceste componente permit modelarea formală a jocului și aplicarea algoritmilor de căutare.

\subsection{Jocuri deterministe}

În jocurile deterministe nu există elemente aleatorii.

Rezultatul unei acțiuni este întotdeauna cunoscut.

Exemple:

\begin{itemize}
\item șah;
\item dame;
\item X și 0;
\item Go.
\end{itemize}

În aceste jocuri, evoluția depinde exclusiv de deciziile jucătorilor.

\subsection{Jocuri stochastice}

În jocurile stochastice intervin evenimente aleatoare.

Rezultatul unei mutări poate depinde de șansă.

Exemple:

\begin{itemize}
\item table;
\item poker;
\item jocurile cu zaruri.
\end{itemize}

Aceste jocuri necesită modele probabilistice suplimentare.

\subsection{Jocuri cu informație completă}

În jocurile cu informație completă fiecare jucător cunoaște întreaga stare a jocului.

Toate mutările și toate informațiile sunt vizibile.

Exemple:

\begin{itemize}
\item șah;
\item dame;
\item X și 0.
\end{itemize}

Aceste jocuri sunt cele mai frecvent utilizate pentru prezentarea algoritmului MiniMax.

\subsection{Jocuri cu informație incompletă}

În jocurile cu informație incompletă anumite informații sunt ascunse.

Jucătorii trebuie să ia decizii fără a cunoaște complet starea jocului.

Exemple:

\begin{itemize}
\item poker;
\item blackjack.
\end{itemize}

În astfel de situații devine necesară estimarea informațiilor lipsă.

\subsection{Jocuri cu doi jucători}

Majoritatea exemplelor studiate în Inteligența Artificială implică doi jucători:

\begin{itemize}
\item MAX;
\item MIN.
\end{itemize}

MAX încearcă să maximizeze câștigul.

MIN încearcă să minimizeze câștigul lui MAX.

Această formulare stă la baza algoritmului MiniMax.

\subsection{Jocuri cu mai mulți jucători}

Există și jocuri care implică mai mult de doi participanți.

Exemple:

\begin{itemize}
\item jocuri de strategie;
\item jocuri de negociere;
\item jocuri cooperative.
\end{itemize}

Acestea necesită modele mai complexe decât cele utilizate în MiniMax.

\subsection{Exemple clasice în Inteligența Artificială}

Cele mai utilizate jocuri pentru testarea algoritmilor de Inteligență Artificială sunt:

\begin{itemize}
\item X și 0;
\item șah;
\item dame;
\item Connect Four;
\item Othello (Reversi);
\item Go.
\end{itemize}

Aceste jocuri oferă un cadru ideal pentru studierea strategiilor de luare a deciziilor.

\subsection{Concluzie}

Pentru a determina mutarea optimă într-un joc este necesară explorarea posibilelor evoluții ale jocului.

Acest lucru conduce la construirea arborelui de joc, care reprezintă fundamentul algoritmului MiniMax.













\section{Jocuri cu sumă zero}

În teoria jocurilor, un caz particular foarte important este reprezentat de jocurile cu sumă zero.

Într-un joc cu sumă zero, câștigul unui jucător este egal cu pierderea adversarului.

Cu alte cuvinte, suma rezultatelor tuturor jucătorilor este întotdeauna egală cu zero.

\subsection{Definiție}

Fie doi jucători:

\begin{itemize}
\item MAX;
\item MIN.
\end{itemize}

Dacă un rezultat oferă lui MAX un câștig de +v, atunci același rezultat produce lui MIN o pierdere de -v.

Formal:

\[
U_{MAX} + U_{MIN} = 0
\]

unde:

\begin{itemize}
\item $U_{MAX}$ reprezintă utilitatea lui MAX;
\item $U_{MIN}$ reprezintă utilitatea lui MIN.
\end{itemize}

\subsection{Exemplu simplu}

Considerăm următoarele rezultate:

\begin{center}
\begin{tabular}{|c|c|}
\hline
Rezultat & Utilitate pentru MAX \\
\hline
Victorie & +1 \\
\hline
Remiză & 0 \\
\hline
Înfrângere & -1 \\
\hline
\end{tabular}
\end{center}

Utilitatea pentru MIN va fi:

\begin{center}
\begin{tabular}{|c|c|}
\hline
Rezultat & Utilitate pentru MIN \\
\hline
Victorie MAX & -1 \\
\hline
Remiză & 0 \\
\hline
Înfrângere MAX & +1 \\
\hline
\end{tabular}
\end{center}

\subsection{Exemple de jocuri cu sumă zero}

Majoritatea jocurilor clasice studiate în Inteligența Artificială sunt jocuri cu sumă zero:

\begin{itemize}
\item șah;
\item dame;
\item X și 0;
\item Connect Four;
\item Reversi;
\item Go.
\end{itemize}

În toate aceste jocuri, avantajul unuia dintre jucători reprezintă dezavantajul adversarului.

\subsection{Ipoteza adversarului perfect}

Algoritmul MiniMax presupune că adversarul joacă perfect.

Astfel:

\begin{itemize}
\item MAX încearcă să maximizeze utilitatea;
\item MIN încearcă să minimizeze utilitatea lui MAX.
\end{itemize}

Fiecare jucător alege întotdeauna mutarea cea mai favorabilă pentru el.

\subsection{Interpretarea unei stări}

O stare a jocului poate fi evaluată printr-o valoare numerică.

Exemplu:

\begin{itemize}
\item +100 = situație foarte favorabilă pentru MAX;
\item 0 = poziție echilibrată;
\item -100 = situație foarte favorabilă pentru MIN.
\end{itemize}

Scopul lui MAX este obținerea celei mai mari valori posibile.

Scopul lui MIN este obținerea celei mai mici valori posibile.

\subsection{De ce sunt importante jocurile cu sumă zero?}

Jocurile cu sumă zero permit formularea matematică simplă a conflictului dintre doi agenți.

Această proprietate face posibilă aplicarea algoritmului MiniMax și a tehnicilor de Alpha-Beta Pruning.

Fără această ipoteză, analiza jocului devine mult mai complexă și necesită modele mai avansate din teoria jocurilor.

\subsection{Concluzie}

În jocurile cu sumă zero, interesul celor doi jucători este complet opus.

MAX încearcă să maximizeze câștigul, iar MIN încearcă să îl minimizeze.

Această idee stă la baza construirii arborelui de joc și a algoritmului MiniMax, care va fi studiat în secțiunile următoare.


















\section{Formularea unei probleme de joc}

Pentru a aplica algoritmi de căutare adversarială, un joc trebuie descris într-o formă matematică precisă.

Această descriere poartă numele de formularea problemei de joc.

Scopul formulării este definirea tuturor elementelor necesare pentru ca un agent să poată analiza situația curentă și să aleagă o mutare optimă.

\subsection{Componentele unei probleme de joc}

O problemă de joc este definită prin:

\begin{itemize}
\item starea inițială;
\item funcția succesor;
\item mulțimea acțiunilor legale;
\item jucătorul curent;
\item testul de terminare;
\item funcția de utilitate.
\end{itemize}

Aceste componente permit construirea arborelui de joc.

\subsection{Starea inițială}

Starea inițială reprezintă configurația de la care începe jocul.

Exemple:

\begin{itemize}
\item tabla de șah la începutul partidei;
\item tabla goală la X și 0;
\item poziția inițială dintr-un joc de dame.
\end{itemize}

Toate celelalte stări se obțin pornind de la această configurație.

\subsection{Acțiunile}

Pentru fiecare stare există o mulțime de acțiuni legale.

O acțiune reprezintă o mutare permisă de regulile jocului.

Exemple:

\begin{itemize}
\item mutarea unui pion;
\item plasarea unui simbol X pe tablă;
\item capturarea unei piese.
\end{itemize}

Numărul acțiunilor poate varia de la o stare la alta.

\subsection{Funcția succesor}

Funcția succesor descrie modul în care se obține o nouă stare după efectuarea unei acțiuni.

Formal:

\[
Succesor(Stare, Actiune)
=
StareNoua
\]

Această funcție permite generarea tuturor mutărilor posibile dintr-o anumită poziție.

\subsection{Jucătorul curent}

În fiecare stare trebuie să știm cine urmează să efectueze mutarea.

În algoritmul MiniMax folosim convențional doi jucători:

\begin{itemize}
\item MAX;
\item MIN.
\end{itemize}

Nivelurile arborelui alternează între MAX și MIN.

\subsection{Stări terminale}

O stare terminală este o stare în care jocul s-a încheiat.

În aceste stări nu mai există acțiuni disponibile.

Exemple:

\begin{itemize}
\item șah mat;
\item remiză;
\item victoria unui jucător la X și 0;
\item ocuparea completă a tablei.
\end{itemize}

\subsection{Testul de terminare}

Testul de terminare determină dacă o stare este terminală.

Formal:

\[
Terminal(Stare)
=
true
\]

dacă jocul s-a încheiat.

În caz contrar:

\[
Terminal(Stare)
=
false
\]

și procesul de căutare continuă.

\subsection{Funcția de utilitate}

Funcția de utilitate asociază o valoare numerică fiecărei stări terminale.

Această valoare exprimă cât de favorabil este rezultatul pentru jucătorul MAX.

Exemplu simplu:

\begin{center}
\begin{tabular}{|c|c|}
\hline
Rezultat & Utilitate \\
\hline
Victorie MAX & +1 \\
\hline
Remiză & 0 \\
\hline
Victorie MIN & -1 \\
\hline
\end{tabular}
\end{center}

Scopul lui MAX este maximizarea utilității.

Scopul lui MIN este minimizarea acesteia.

\subsection{Exemplu: X și 0}

O formulare simplificată pentru jocul X și 0 este:

\begin{itemize}
\item starea inițială: tabla goală;
\item acțiuni: plasarea unui simbol într-o celulă liberă;
\item funcția succesor: actualizarea tablei;
\item test terminal: trei simboluri identice pe linie, coloană sau diagonală ori tablă completă;
\item utilitate:
    \begin{itemize}
    \item +1 pentru victoria lui MAX;
    \item 0 pentru remiză;
    \item -1 pentru victoria lui MIN.
    \end{itemize}
\end{itemize}

\subsection{Concluzie}

După formularea problemei putem genera toate evoluțiile posibile ale jocului.

Aceste evoluții sunt reprezentate prin intermediul arborelui de joc, care constituie baza algoritmului MiniMax.















\section{Arborele de joc}

Pentru a determina cea mai bună mutare într-un joc, este necesară explorarea tuturor evoluțiilor posibile ale acestuia.

Aceste evoluții sunt reprezentate prin intermediul arborelui de joc.

Arborele de joc reprezintă o structură arborescentă care descrie toate mutările posibile ce pot fi efectuate pornind dintr-o anumită stare.

\subsection{Definiție}

Un arbore de joc este un arbore în care:

\begin{itemize}
\item rădăcina reprezintă starea curentă;
\item fiecare muchie reprezintă o mutare;
\item fiecare nod reprezintă o stare a jocului;
\item frunzele reprezintă stări terminale sau stări evaluate.
\end{itemize}

Scopul algoritmilor de căutare este determinarea celui mai bun drum prin arbore.

\subsection{Structura arborelui}

Într-un arbore de joc distingem:

\begin{itemize}
\item nodul rădăcină;
\item noduri interne;
\item noduri frunză;
\item muchii.
\end{itemize}

Fiecare nod intern poate genera unul sau mai mulți succesori.

Numărul de succesori reprezintă factorul de ramificare al arborelui.

\subsection{Alternarea jucătorilor}

Nivelurile arborelui alternează între cei doi jucători.

Astfel:

\begin{itemize}
\item nivelul 0 aparține lui MAX;
\item nivelul 1 aparține lui MIN;
\item nivelul 2 aparține lui MAX;
\item nivelul 3 aparține lui MIN;
\item și așa mai departe.
\end{itemize}

Fiecare jucător alege mutarea care îi maximizează propriul avantaj.

\subsection{Exemplu simplu}

Considerăm următorul arbore:

\begin{verbatim}
                                 MAX
                    /-------------|-------------\
                  MIN            MIN           MIN
                /-----\        /-----\       /-----\
              MAX    MAX     MAX    MAX    MAX     MAX
             /  \    /  \   /  \   /  \   /  \    /  \
            3    5  6    9 1    2 0    7 4    8  2    6
\end{verbatim}

Rădăcina reprezintă starea curentă.

MAX trebuie să aleagă una dintre mutările disponibile.

Pentru a lua decizia corectă este necesar să analizăm valorile care se propagă din frunze către rădăcină.

\subsection{Frunzele arborelui}

Frunzele sunt stările finale la care se poate ajunge.

În aceste noduri nu mai există mutări disponibile.

Fiecărei frunze i se asociază o valoare numerică denumită utilitate.

Exemplu:

\begin{itemize}
\item +100 pentru victorie;
\item 0 pentru remiză;
\item -100 pentru înfrângere.
\end{itemize}

Aceste valori permit compararea diferitelor linii de joc.

\subsection{Adâncimea arborelui}

Adâncimea reprezintă numărul maxim de mutări analizate de la rădăcină până la o frunză.

Cu cât adâncimea este mai mare:

\begin{itemize}
\item se analizează mai multe poziții;
\item crește acuratețea deciziei;
\item crește timpul necesar căutării.
\end{itemize}

În jocuri complexe precum șahul, explorarea completă a arborelui este imposibilă.

Din acest motiv sunt utilizate limite de adâncime și funcții de evaluare.

\subsection{Factorul de ramificare}

Factorul de ramificare reprezintă numărul mediu de mutări posibile dintr-o stare.

Exemple aproximative:

\begin{itemize}
\item X și 0: până la 9 mutări;
\item dame: aproximativ 8 mutări;
\item șah: aproximativ 35 mutări;
\item Go: peste 200 mutări.
\end{itemize}

Un factor de ramificare mare conduce la creșterea rapidă a dimensiunii arborelui.

\subsection{Generarea arborelui}

Pornind de la starea curentă:

\begin{enumerate}
\item se generează toate acțiunile posibile;
\item pentru fiecare acțiune se obține o stare nouă;
\item procesul se repetă pentru fiecare succesor;
\item se continuă până la atingerea unei stări terminale sau a unei limite de adâncime.
\end{enumerate}

Rezultatul este arborele de joc.

\subsection{Importanța arborelui de joc}

Arborele de joc reprezintă baza algoritmului MiniMax.

Acesta permite simularea comportamentului celor doi jucători și determinarea mutării optime.

Fără arborele de joc nu putem aplica algoritmi de căutare adversarială.

\subsection{Concluzie}

Arborele de joc descrie toate evoluțiile posibile ale unei partide.

Fiecare nod reprezintă o stare, iar fiecare muchie reprezintă o mutare.

Analizarea acestor arbori permite alegerea celei mai bune decizii prin intermediul algoritmului MiniMax.
























\section{Valoarea unei stări}

În cadrul căutării adversariale, fiecare stare a jocului trebuie evaluată pentru a determina cât de favorabilă este pentru fiecare jucător.

Această evaluare poartă numele de valoarea unei stări.

Valoarea unei stări reprezintă o estimare numerică a avantajului pe care îl are un jucător într-o anumită poziție.

\subsection{Scopul evaluării}

Atunci când un agent trebuie să aleagă o mutare, acesta compară valorile stărilor care pot fi obținute.

În general:

\begin{itemize}
\item valorile mari sunt favorabile pentru MAX;
\item valorile mici sunt favorabile pentru MIN.
\end{itemize}

Prin urmare:

\begin{itemize}
\item MAX încearcă să maximizeze valoarea;
\item MIN încearcă să minimizeze valoarea.
\end{itemize}

\subsection{Valori terminale}

Pentru stările terminale, valoarea este cunoscută exact.

Exemplu:

\begin{center}
\begin{tabular}{|c|c|}
\hline
Rezultat & Valoare \\
\hline
Victorie MAX & +1 \\
\hline
Remiză & 0 \\
\hline
Victorie MIN & -1 \\
\hline
\end{tabular}
\end{center}

Aceste valori reprezintă utilitatea finală a jocului.

\subsection{Valori pentru stări intermediare}

Majoritatea stărilor întâlnite în arbore nu sunt terminale.

În aceste situații trebuie determinată valoarea stării pe baza mutărilor care pot urma.

Algoritmul MiniMax calculează această valoare propagând informația de la frunze către rădăcină.

\subsection{Noduri MAX}

Într-un nod MAX, jucătorul poate alege dintre mai multe mutări.

Se presupune că va selecta mutarea cea mai avantajoasă.

Prin urmare:

\[
Val(MAX)=\max(Val(Succesori))
\]

Exemplu:

\begin{verbatim}
      MAX
     / | \
    3  5  2
\end{verbatim}

Valoarea nodului este:

\[
Val(MAX)=5
\]

deoarece MAX va prefera valoarea cea mai mare.

\subsection{Noduri MIN}

Într-un nod MIN, adversarul încearcă să reducă avantajul lui MAX.

Prin urmare:

\[
Val(MIN)=\min(Val(Succesori))
\]

Exemplu:

\begin{verbatim}
      MIN
     / | \
    3  5  2
\end{verbatim}

Valoarea nodului este:

\[
Val(MIN)=2
\]

deoarece MIN va selecta valoarea cea mai mică.

\subsection{Exemplu complet}

Considerăm arborele:

\begin{verbatim}
               MAX
            /       \
          MIN       MIN
         /  \       /  \
        3    5     2    9
\end{verbatim}

Calculăm:

\[
Val(MIN_{st})=\min(3,5)=3
\]

\[
Val(MIN_{dr})=\min(2,9)=2
\]

Apoi:

\[
Val(MAX)=\max(3,2)=3
\]

Prin urmare, mutarea optimă pentru MAX este cea care conduce la valoarea 3.

\subsection{Interpretarea valorilor}

Semnul valorii este mai important decât valoarea absolută.

În general:

\begin{itemize}
\item valoare pozitivă → avantaj pentru MAX;
\item valoare negativă → avantaj pentru MIN;
\item valoare apropiată de 0 → poziție echilibrată.
\end{itemize}

Exemplu:

\begin{itemize}
\item +100 poate indica o victorie aproape sigură pentru MAX;
\item -100 poate indica o victorie aproape sigură pentru MIN;
\item +5 poate reprezenta un avantaj redus pentru MAX.
\end{itemize}

\subsection{Propagarea valorilor}

Ideea fundamentală a algoritmului MiniMax este propagarea valorilor de la frunzele arborelui către rădăcină.

Procesul este:

\begin{enumerate}
\item se evaluează frunzele;
\item se calculează valorile nodurilor MIN;
\item se calculează valorile nodurilor MAX;
\item se ajunge la valoarea rădăcinii.
\end{enumerate}

Valoarea obținută în rădăcină reprezintă rezultatul pe care MAX îl poate garanta dacă ambii jucători joacă optim.

\subsection{Concluzie}

Valoarea unei stări reprezintă fundamentul algoritmului MiniMax.

Fiecare jucător presupune că adversarul va lua întotdeauna cea mai bună decizie pentru propriul interes, iar valoarea unui nod este calculată folosind operațiile de maxim și minim asupra succesorilor săi.






\enddocument{}